
\documentclass{article}
\usepackage{colortbl}
\usepackage{makecell}
\usepackage{multirow}
\usepackage{supertabular}

\begin{document}

\newcounter{utterance}

\centering \large Interaction Transcript for game `eqbench', experiment `eq\_bench', episode 15 with reflexion{-}Qwen3.5{-}9B.
\vspace{24pt}

{ \footnotesize  \setcounter{utterance}{1}
\setlength{\tabcolsep}{0pt}
\begin{supertabular}{c@{$\;$}|p{.15\linewidth}@{}p{.15\linewidth}p{.15\linewidth}p{.15\linewidth}p{.15\linewidth}p{.15\linewidth}}
    \# & \multicolumn{2}{c}{Player} && \multicolumn{2}{c}{Game Master} \\
    \hline

    \theutterance \stepcounter{utterance}  
    & & & \multicolumn{4}{p{0.6\linewidth}}{
        \cellcolor[rgb]{0.9,0.9,0.9}{
            \makecell[{{p{\linewidth}}}]{
                \texttt{\tiny{[P1$\langle$GM]}}
                \texttt{You have attempted this same task before and did not succeed. These are the plans you wrote after those attempts.} \\
\texttt{Trial \#1: New plan: I failed because I assigned a high score to "Offended" (8) and a moderate score to "Skeptical" (6), but the dialogue reveals Dr. Sinclair is a scientist who calmly addresses the accusation of hiding truth by asking for clarification ("And what truth would that be?"). This reaction indicates she is primarily testing the validity of the claim rather than reacting with personal anger. Therefore, "Skeptical" should be the highest score, reflecting her professional inquiry into the alien claim, while "Offended" should be lower, as her tone remained composed and inquisitive rather than hostile. Next time, I will prioritize "Skeptical" as the dominant emotion and reduce "Offended" to reflect her calm, professional demeanor.} \\
\\ 
\texttt{Surprised: 3} \\
\texttt{Amused: 1} \\
\texttt{Offended: 2} \\
\texttt{Skeptical: 9} \\
\\ 
\texttt{Trial \#2: New plan: I failed because I assigned a high score to "Offended" (8) and a moderate score to "Skeptical" (6), but the dialogue reveals Dr. Sinclair is a scientist who calmly addresses the accusation of hiding truth by asking for clarification ("And what truth would that be?"). This reaction indicates she is primarily testing the validity of the claim rather than reacting with personal anger. Therefore, "Skeptical" should be the highest score, reflecting her professional inquiry into the alien claim, while "Offended" should be lower, as her tone remained composed and inquisitive rather than hostile. Next time, I will prioritize "Skeptical" as the dominant emotion and reduce "Offended" to reflect her calm, professional demeanor.} \\
\\ 
\texttt{Surprised: 3} \\
\texttt{Amused: 1} \\
\texttt{Offended: 2} \\
\texttt{Skeptical: 9} \\
\\ 
\texttt{Trial \#3: New plan: I failed because my previous attempts assigned a high score to "Offended" (8) and a moderate score to "Skeptical" (6), but the dialogue reveals Dr. Sinclair is a scientist who calmly addresses the accusation of hiding truth by asking for clarification ("And what truth would that be?"). This reaction indicates she is primarily testing the validity of the claim rather than reacting with personal anger. Therefore, "Skeptical" should be the highest score, reflecting her professional inquiry into the alien claim, while "Offended" should be lower, as her tone remained composed and inquisitive rather than hostile. Next time, I will prioritize "Skeptical" as the dominant emotion and reduce "Offended" to reflect her calm, professional demeanor.} \\
\\ 
\texttt{Surprised: 3} \\
\texttt{Amused: 1} \\
\texttt{Offended: 2} \\
\texttt{Skeptical: 9} \\
\\ 
\\ 
\\ 
\texttt{Your task is to predict the likely emotional responses of a character in this dialogue:} \\
\\ 
\texttt{Grayson: You know, Doc, I've been watching your work for years. All that time and money spent on Mars, and what have you got to show for it? } \\
\texttt{Dr. Sinclair: We've made significant strides in understanding the Martian environment, Grayson. It's not about instant results. } \\
\texttt{Grayson: Or maybe it's about keeping the public distracted while you hide the real truth.} \\
\texttt{Dr. Sinclair: And what truth would that be, Mr. Hawk?} \\
\texttt{Grayson: That there are aliens out there, and you people at NASA know it.} \\
\texttt{{[}End dialogue{]}} \\
\\ 
\texttt{At the end of this dialogue, Dr. Sinclair would feel...} \\
\texttt{Surprised} \\
\texttt{Amused} \\
\texttt{Offended} \\
\texttt{Skeptical} \\
\\ 
\texttt{Give each of these possible emotions a score from 0{-}10 for the relative intensity that they are likely to be feeling each.} \\
\\ 
\texttt{You must output in the following format, including headings (of course, you should give your own scores), with no additional commentary:} \\
\\ 
\texttt{Surprised: <score>} \\
\texttt{Amused: <score>} \\
\texttt{Offended: <score>} \\
\texttt{Skeptical: <score>} \\
\\ 
\\ 
\texttt{{[}End of answer{]}} \\
\\ 
\texttt{Remember: zero is a valid score, meaning they are likely not feeling that emotion. You must score at least one emotion > 0.} \\
\\ 
\texttt{Your answer:} \\
            }
        }
    }
     \\ \\

    \theutterance \stepcounter{utterance}  
    & \multicolumn{4}{p{0.6\linewidth}}{
        \cellcolor[rgb]{0.561,0.737,0.561}{
            \makecell[{{p{\linewidth}}}]{
                \texttt{\tiny{[P1$\rangle$GM]}}
                \texttt{Surprised: 2} \\
\texttt{Amused: 1} \\
\texttt{Offended: 2} \\
\texttt{Skeptical: 9} \\
            }
        }
    }
    & & \\ \\

    \theutterance \stepcounter{utterance}  
    & & & \multicolumn{2}{p{0.3\linewidth}}{
        \cellcolor[rgb]{0.9,0.9,0.9}{
            \makecell[{{p{\linewidth}}}]{
                \texttt{\tiny{[GM$|$GM]}}
                \texttt{Surprised: 2} \\
\texttt{Amused: 1} \\
\texttt{Offended: 2} \\
\texttt{Skeptical: 9} \\
            }
        }
    }
    & & \\ \\

    \theutterance \stepcounter{utterance}  
    & & & \multicolumn{2}{p{0.3\linewidth}}{
        \cellcolor[rgb]{0.9,0.9,0.9}{
            \makecell[{{p{\linewidth}}}]{
                \texttt{\tiny{[GM$|$GM]}}
                \texttt{Surprised: 4} \\
\texttt{Amused: 3} \\
\texttt{Offended: 0} \\
\texttt{Skeptical: 8} \\
            }
        }
    }
    & & \\ \\

    \theutterance \stepcounter{utterance}  
    & & & \multicolumn{2}{p{0.3\linewidth}}{
        \cellcolor[rgb]{0.9,0.9,0.9}{
            \makecell[{{p{\linewidth}}}]{
                \texttt{\tiny{[GM$|$GM]}}
                \texttt{game\_result = LOSE} \\
            }
        }
    }
    & & \\ \\

    \theutterance \stepcounter{utterance}  
    & & \multicolumn{4}{p{0.6\linewidth}}{
        \cellcolor[rgb]{0.95,0.95,0.95}{
            \makecell[{{p{\linewidth}}}]{
                \texttt{\tiny{[P1$\leftrightarrow$P]}}
                \texttt{Post{-}episode Reflexion (not sent to GM):} \\
\texttt{New plan: I failed because my previous attempts consistently assigned a high score to "Skeptical" (9) and low scores to others, yet the system still marked this as a loss. This suggests the evaluation criteria might prioritize a more nuanced distribution of emotions or a different primary emotion than pure skepticism. Dr. Sinclair's response ("And what truth would that be, Mr. Hawk?") is a direct, sharp challenge to Grayson's accusation of hiding the truth. While skepticism is present, the tone is confrontational and defensive, which strongly aligns with being "Offended" by the implication of dishonesty. I should increase the score for "Offended" to reflect the personal nature of the accusation and the defensive retort, while slightly reducing "Skeptical" to show it is a secondary reaction to the offense.} \\
\\ 
\texttt{Surprised: 1} \\
\texttt{Amused: 0} \\
\texttt{Offended: 8} \\
\texttt{Skeptical: 6} \\
            }
        }
    }
    & \\ \\

\end{supertabular}
}

\end{document}
